The thesis develops a distributional analysis for Richardson--Romberg
averaged constant-stepsize linear stochastic approximation under Markovian
noise. The work starts from the usual LSA recursion and isolates the effect
of coupling two step sizes, $\alpha$ and $2\alpha$, in the Polyak--Ruppert
average. The main mathematical contribution is a proof strategy that combines
deterministic RR weight bounds, Poisson martingale reduction for Markovian
noise, predictable-variance comparison, and explicit control of the
misadjustment terms.

The stationary part of the analysis gives a Berry--Esseen assembly for scalar
RR statistics under the augmented-chain convention. At the balanced scale
$\alpha=c n^{-1/2}$, the result identifies the covariance target
$\Sigma_\infty=\overline A^{-1}\Sigma_{\varepsilon}^{(M)}\overline A^{-\top}$
and gives a non-asymptotic normal approximation up to logarithmic factors.
The deterministic-start part then transfers the stationary theorem to the
practical burned-in estimator by controlling deterministic transients,
initial random products, startup discrepancies, and finite-window
normalization errors.

The numerical experiments are consistent with the theoretical picture.
Step-size Richardson--Romberg extrapolation mainly improves the center of the
confidence interval by reducing constant-stepsize bias, whereas OBM and
lugsail OBM affect the long-run variance estimate. The experiments also show
that lugsail corrections should be interpreted as variance-estimator
bias-reduction tools whose effect depends on the block size and dependence
regime.

Further work should complete three directions. First, the current scalar
Berry--Esseen theory should be extended to more explicitly multivariate
confidence regions. Second, the OBM and lugsail covariance estimators should
be analyzed non-asymptotically along RR-averaged constant-stepsize LSA
trajectories. Third, the dependence on the mixing time and the Lyapunov
contraction constants should be sharpened so that the theory gives more
direct practical guidance for choosing burn-in lengths, step sizes, and block
sizes.

